In this final project, the Metropolis and HMC algorithms were implemented for the two-dimensional $\phi^4$-theory and the performance was compared across different integrators, programming languages and physical regimes. 
\\\\
Before performing the main anaylsis, the algorithms were tuned to the correct acceptance rates, see Section \ref{sec: tuning}. The Metropolis algorithm was tuned to approximately 50\% acceptance rate, HMC was tuned to around 90\%. 
\\\\
In Section \ref{sec:chain_visual}, the Monte Carlo history of the mean field observable was visualized to assess thermalization. These plots were found not be of great quantitative usefulness. However, it could be concluded that the Metropolis algorithm has the longest thermalization time, since it has local updates, in contrast to the global molecular-dynamics HMC updates. 
\\\\
Section \ref{sec: data_runtime} focused on the run times of the algorithms (Metropolis vs. HMC with both integrators) implemented in two programming languages (Python vs. Python). Translating the code from Python to Rust yielded a substantial speedup of a factor $4-6$. It was conlcluded that implementing expensive tasks in Rust is beneficial. The use of PyO3 enables a convenient framework in which performance-critical components are written in Rust, while data analysis and visualization remain in Python. In general, the Metropolis algorithm is computationally inexpensive. Solving the equations of motions for HMC numerically is much more expensive, the use of integrator does not influence the run time of HMC dramatically. Additionally, it was observed that the run time depends on the phase of the system, which can be interpreted in terms of critical slowing down. 
\\\\
In Section \ref{sec: errs}, statistical errors and integrated autocorrelation times were analyzed. The results show that the Metropolis algorithm produces significantly larger statistical errors than HMC. Both leapfrog and omelyan2 integrators exhibit similar statistical uncertainties, but HMC clearly outperforms Metropolis in terms of autocorrelation times. The statistical errors are larger in the broken phase than in the symmetric phase, which can be attributed to the presence of two degenerate vacua. Furthermore, while statistical errors decrease with increasing lattice size $L$, the integrated autocorrelation time increases, reflecting critical slowing down. This effect is particularly pronounced for Metropolis, whereas HMC suppresses it efficiently. Discrepancies between Python and Rust implementations are not fully explained, this needs further investigations, since errors and autocorrelation times should be independent of the programming language. 
\\\\
Section \ref{sec: int_comparison} examines the performance of the integrators leapfrog and omelyan2 in HMC. The HMC algorithm, with both leapfrog and omelyan2 integrators, demonstrates excellent energy conservation and high acceptance rates, particularly with the omelyan2 integrator. Our results show that omelyan2 consistently produces smaller Hamiltonian violations than leapfrog, confirming its improved stability and robustness across the symmetric and broken phases. The step size $\varepsilon$ strongly influences energy conservation, with the standard deviation of $\Delta H$ growing as approximately $\varepsilon^2$.
\\\\
The study of the topological charge $Q_R$ using the Rust Metropolis implementation further illustrates the physical differences between the phases. In the broken phase, non-trivial topological excitations such as kinks and anti-kinks emerge, with $Q_R$ fluctuating around non-zero values and occasionally changing sign due to tunneling events. In the symmetric phase, these topological effects are suppressed, and $Q_R$ remains close to zero, as expected.
\\\\
Overall, our analysis confirms that the choice of algorithm, integrator, and programming language has a strong impact on simulation efficiency, statistical precision, and energy conservation. HMC with omelyan2 integrator in Rust provides the most favorable combination of accuracy, efficiency, and low autocorrelation.
\section{Possible Extensions}
Here are a few suggestions on how to extend this final project: 
\begin{itemize}
    \item The current code base of this project is written such that the action of the theory could be exchanged with relatively simply, allowing the code to be used for different, more realistic theories than scalar $\phi^4$ theory. However, it would require some effort to implement fermions or gauge fields. 
    \item While the leapfrog and especially omelyan2 integrators for HMC are robust and stable, higher order integrators could be added in order to further reduce integration errors and to improve energy conservation. 
    \item It would be interesting to study more sophisticated observables. Beyond the mean field and the topological charge, one could measure for example correclation functions or susceptibilities. 
    \item The Rust implementation could be improved significantly using parallelization or GPU acceleration, allowing simulations on larger lattices or large parameter scans. Combining high-performance Rust code with Python plotting and analysis by means of PyO3 provides many possibilities for LFT. 
\end{itemize}